% ============================================================
%  preamble.tex — packages and global configuration.
%  Derived from marlinjai/academic-thesis-template, adapted for
%  a Master's thesis and split into modular files.
% ============================================================

% --- Codifica e lingue (italiano principale; inglese e tedesco disponibili
%     per citazioni e per un eventuale abstract in seconda lingua).
%     In babel l'ultima lingua elencata e' quella principale.
\usepackage[utf8]{inputenc}
\usepackage[T1]{fontenc}
\usepackage[english,ngerman,italian]{babel}
\usepackage[autostyle,italian=guillemets]{csquotes}
\usepackage{lmodern}
\usepackage{microtype}

% --- Math
\usepackage{amsmath, amssymb, amsthm}

% --- Graphics and figures
\usepackage{graphicx}
\graphicspath{{figures/}}
\usepackage{float}

% --- Tables and captions
\usepackage{booktabs}
\usepackage{longtable}
\usepackage{tabularx}
\usepackage{threeparttable}   % table notes under regression tables
\usepackage{multirow}
\usepackage[small, labelfont=bf]{caption}
\usepackage{subcaption}

\usepackage{eurosym}

% --- Numbers and units
\usepackage{siunitx}
\sisetup{group-separator={,}, group-minimum-digits=4, detect-weight=true}

% --- Page layout
\usepackage[hmargin=2.5cm, top=3cm, bottom=3cm]{geometry}

% --- Bibliografia (autore-anno, biblatex + biber)
% Convenzione economica. Per una tesi di area informatica sostituire con:
%   \usepackage[backend=biber, style=ieee, sorting=none]{biblatex}
% e usare \cite{} al posto di \textcite{}/\parencite{}.
\usepackage[backend=biber, style=authoryear, language=autobib,
            maxcitenames=2, maxbibnames=99,
            uniquename=false, uniquelist=false, natbib=true, dashed=false,
            giveninits=true]{biblatex}
\addbibresource{references.bib}

% --- Glossary and acronyms
\usepackage[acronym,toc]{glossaries}
\makeglossaries
% ============================================================
%  glossary.tex: Glossareintraege und Abkuerzungen.
%  Aufruf mit \gls{...}, \glspl{...}, \Gls{...},
%  oder \acrshort{...} / \acrlong{...} / \acrfull{...}.
%  Gedruckt werden nur tatsaechlich zitierte Eintraege.
% ============================================================

\newglossaryentry{altholz}{
  name=Altholz,
  description={Gebrauchtes Holz aus Rueckbau, Abbruch oder Demontage, das
               nach der Altholzverordnung in vier Kategorien eingeteilt wird;
               nur die gering belasteten Kategorien kommen fuer eine tragende
               Wiederverwendung in Betracht}
}

\newglossaryentry{druckzone}{
  name=Druckzone,
  description={Der Querschnittsbereich eines auf Biegung beanspruchten
               Traegers oberhalb der neutralen Faser, in dem die
               Laengsspannungen Druckspannungen sind}
}

\newglossaryentry{zugzone}{
  name=Zugzone,
  description={Der Querschnittsbereich eines auf Biegung beanspruchten
               Traegers unterhalb der neutralen Faser, in dem die
               Laengsspannungen Zugspannungen sind}
}

\newglossaryentry{neutralefaser}{
  name=neutrale Faser,
  description={Die spannungsfreie Schicht im Querschnitt eines gebogenen
               Traegers, die Druck- und Zugzone voneinander trennt}
}

\newglossaryentry{lastgeschichte}{
  name=Lastgeschichte,
  description={Die gesamte Abfolge der Beanspruchungen, die ein Bauteil
               waehrend seiner Nutzungsdauer erfahren hat, einschliesslich
               Dauer, Hoehe und Wechsel der Lasten sowie der begleitenden
               Feuchteverhaeltnisse}
}

\newglossaryentry{kriechen}{
  name=Kriechen,
  description={Die zeitabhaengige, mit der Belastungsdauer zunehmende
               Verformung unter konstanter Last; bei Holz durch Feuchtewechsel
               erheblich verstaerkt (mechano-sorptives Kriechen)}
}

\newglossaryentry{restfestigkeit}{
  name=Restfestigkeit,
  description={Die nach einer Nutzungsphase verbliebene Tragfaehigkeit eines
               Bauteils, im Unterschied zur Festigkeit des ungenutzten
               Ausgangsmaterials}
}

\newglossaryentry{rohdichte}{
  name=Rohdichte,
  description={Masse je Volumeneinheit des Holzes bei definierter Holzfeuchte;
               wichtigster Einzelindikator fuer die mechanischen Eigenschaften
               und daher als Kovariate im Zonenvergleich zu fuehren}
}

\newglossaryentry{sortierung}{
  name=Sortierung,
  description={Die Zuordnung von Bauholz zu einer Festigkeitsklasse, entweder
               visuell nach DIN 4074 oder maschinell nach EN 14081}
}

% --- Abkuerzungen
\newacronym{altholzv}{AltholzV}{Altholzverordnung}
\newacronym{krwg}{KrWG}{Kreislaufwirtschaftsgesetz}
\newacronym{hnee}{HNEE}{Hochschule fuer nachhaltige Entwicklung Eberswalde}
\newacronym{mor}{MOR}{Biegefestigkeit (modulus of rupture)}
\newacronym{moe}{MOE}{Elastizitaetsmodul (modulus of elasticity)}
\newacronym{ndt}{NDT}{zerstoerungsfreie Pruefung (non-destructive testing)}
\newacronym{dol}{DOL}{Lastdauereffekt (duration of load)}
\newacronym{anova}{ANOVA}{Varianzanalyse (analysis of variance)}


% --- Code listings
\usepackage{listings}
\usepackage{xcolor}
\definecolor{codegray}{gray}{0.95}
\definecolor{codekeyword}{RGB}{0,0,180}
\definecolor{codecomment}{RGB}{0,128,0}
\definecolor{codestring}{RGB}{163,21,21}
\lstset{
  basicstyle=\footnotesize\ttfamily,
  backgroundcolor=\color{codegray},
  keywordstyle=\color{codekeyword}\bfseries,
  commentstyle=\color{codecomment}\itshape,
  stringstyle=\color{codestring},
  frame=single,
  breaklines=true,
  breakatwhitespace=true,
  showstringspaces=false,
  numbers=left,
  numberstyle=\tiny,
  captionpos=b,
  tabsize=2,
  language=Python % default; override per listing with [language=...]
}

% Denominazioni italiane per i listati di codice.
\renewcommand{\lstlistlistingname}{Elenco dei listati}
\renewcommand{\lstlistingname}{Listato}

% --- Hyperlinks (load late, as hyperref likes to be near-last)
% hypertexnames=false: the appendix restarts page numbering, which otherwise
% produces "duplicate destination" warnings on every build.
\usepackage[hidelinks,plainpages=false,pdfpagelabels,hypertexnames=false]{hyperref}
\usepackage{cleveref}
% Denominazioni italiane esplicite: non dipendono dal rilevamento automatico
% della lingua da parte di cleveref.
\crefname{chapter}{capitolo}{capitoli}
\Crefname{chapter}{Capitolo}{Capitoli}
\crefname{section}{sezione}{sezioni}
\Crefname{section}{Sezione}{Sezioni}
\crefname{subsection}{sezione}{sezioni}
\Crefname{subsection}{Sezione}{Sezioni}
\crefname{equation}{equazione}{equazioni}
\Crefname{equation}{Equazione}{Equazioni}
\crefname{figure}{figura}{figure}
\Crefname{figure}{Figura}{Figure}
\crefname{table}{tabella}{tabelle}
\Crefname{table}{Tabella}{Tabelle}
\crefname{listing}{listato}{listati}
\Crefname{listing}{Listato}{Listati}

% --- Ambienti di tipo teorema (denominazioni italiane)
\newtheorem{theorem}{Teorema}[chapter]
\newtheorem{definition}[theorem]{Definizione}
\newtheorem{lemma}[theorem]{Lemma}
\newtheorem{remark}[theorem]{Osservazione}
\newtheorem{corollary}[theorem]{Corollario}
\newtheorem{example}[theorem]{Esempio}
